%=====================================================================
%  Capstone I -- Universidad ESAN -- Ingenieria de Inteligencia Artificial
%  Proyecto: Sistema de Computer Vision para estimacion de aglomeraciones
%            de personas en el Aeropuerto Internacional Jorge Chavez
%  Cliente:  Lima Airport Partners (LAP)
%
%  Compilacion: pdfLaTeX + Biber + pdfLaTeX + pdfLaTeX
%               latexmk -pdf main.tex
%
%  Mantenimiento del informe: Fabian Moreno Ugarte
%  (Validacion Etica/Legal y Documentacion)
%
%  Los archivos .tex estan codificados en UTF-8 sin BOM.
%=====================================================================
\documentclass[11pt,a4paper]{article}

%--------------------- Codificacion e idioma -------------------------
\usepackage[T1]{fontenc}
\usepackage[utf8]{inputenc}
\usepackage{lmodern}
% es-tabla     -> "Tabla" en lugar de "Cuadro"
% es-noquoting -> evita que babel-spanish capture el caracter " , que
%                 entra en conflicto con csquotes/biblatex
\usepackage[spanish,es-tabla,es-noquoting,es-lcroman]{babel}

%--------------------- Geometria y estilo ----------------------------
\usepackage[a4paper,left=3cm,right=2.5cm,top=3cm,bottom=2.5cm]{geometry}
\usepackage{fancyhdr}
\usepackage{setspace}
\usepackage{parskip}

%--------------------- Tablas, listas, graficos ----------------------
\usepackage{graphicx}
\usepackage{xcolor}
\usepackage{booktabs}
\usepackage{array}
\usepackage{tabularx}
\usepackage{enumitem}
\usepackage{ragged2e}

% Tipos de columna propios para las tablas del informe.
% En columnas estrechas el texto justificado produce lineas muy sueltas
% (avisos "Underfull \hbox"); se usa alineacion a la izquierda, que ademas
% es la opcion tipografica correcta a estos anchos.
%   L{ancho} -> columna de ancho fijo, alineada a la izquierda
%   Y        -> como la columna X de tabularx, alineada a la izquierda
\newcolumntype{L}[1]{>{\RaggedRight\arraybackslash}p{#1}}
\newcolumntype{Y}{>{\RaggedRight\arraybackslash}X}

% NOTA DE COMPATIBILIDAD (longtable)
% Las versiones v4.24 a v4.26 de longtable (oct-dic 2025) tienen una
% regresion: al partir entre paginas una tabla con columnas de tipo p{},
% pdfTeX emite
%     "Infinite glue shrinkage found in box being split"
% y devuelve codigo de salida distinto de cero, aunque igualmente genere
% el PDF. No depende de este documento: se reproduce con un longtable
% minimo de una sola columna p{} que cruce un salto de pagina.
% Esta corregido a partir de v4.27 (2026-01-28), que es la version con la
% que se verifico este proyecto.
% Si en Overleaf aparece ese error, hay dos salidas equivalentes:
%   (a) cambiar la version de TeX Live del proyecto (Menu > Compilador)
%       a una que traiga longtable <= v4.23 o >= v4.27, o
%   (b) pedir la version estable antigua cambiando la linea siguiente por
%       \usepackage{longtable}[=v4.13]
%       (funciona, pero avisa "Command \LT@p@ftntext has changed").
\usepackage{longtable}

%--------------------- Bibliografia ----------------------------------
% csquotes DEBE cargarse antes que biblatex.
\usepackage{csquotes}
% Estilo IEEE numerico: llamadas [1], [4] numeradas por ORDEN DE APARICION.
% sorting=none es lo que garantiza ese orden (no alfabetico); es ademas el
% criterio propio de IEEE, de modo que la opcion es redundante pero se deja
% explicita para que no se pierda si alguien cambia el estilo.
% El estilo anterior era numeric-comp, que ya numeraba igual; lo que cambia
% con ieee es el FORMATO de las entradas y que las citas multiples se
% escriben [2], [3] en vez de [2, 3], como corresponde a IEEE.
\usepackage[
  backend=biber,
  style=ieee,
  sorting=none,
  maxbibnames=99,
  giveninits=true,
  isbn=false,
  doi=false,
  url=false
]{biblatex}
\addbibresource{referencias.bib}

% biblatex-ieee agrupa los millares tambien en los rangos de paginas, lo que
% imprime "pags. 29 363-29 373" en la entrada [3]. En paginacion el separador
% de millares no corresponde, asi que se restaura el formato plano.
\DeclareFieldFormat{pages}{\mkpageprefix[bookpagination]{#1}}

%--------------------- Colores y macros propias ----------------------
\definecolor{azulesan}{RGB}{0,51,102}
\definecolor{grisclaro}{RGB}{242,242,242}

% Caja de aviso reutilizable, construida solo con \fcolorbox + minipage
% para no depender de tcolorbox ni mdframed.
\newcommand{\cajaaviso}[2]{%
  \par\vspace{0.6em}\noindent
  \fcolorbox{azulesan}{grisclaro}{%
    \begin{minipage}{\dimexpr\linewidth-2\fboxsep-2\fboxrule\relax}
      \vspace{0.3em}
      \textbf{\color{azulesan}#1}\par\vspace{0.35em}
      \small #2
      \vspace{0.3em}
    \end{minipage}%
  }\par\vspace{0.6em}
}

% NOTA SOBRE LOS MARCADORES RETIRADOS
% Este documento tuvo una macro \pendiente{tipo}{texto} que imprimia en
% ambar, dentro del cuerpo, los puntos que requerian decision o
% verificacion. El cuerpo ya no lleva ninguno: todos se trasladaron a
% PENDIENTES.md, agrupados por quien puede resolverlos.
% La macro se retiro a proposito y no se dejo definida como no-op: si
% alguien vuelve a escribir \pendiente{...}{...} en una seccion, la
% compilacion falla de inmediato. Esa es la salvaguarda que impide que un
% marcador reaparezca en el PDF sin que nadie lo advierta.
%
% MARCADORES [VERIFICAR] -- MECANISMO (2026-09-13)
% \pendiente y \verificar hacen cosas OPUESTAS y por eso conviven sin
% contradecirse:
%   \pendiente  no existe  -> si reaparece, la compilacion FALLA.
%   \verificar  existe     -> el texto queda en el fuente y NO se imprime.
% \verificar{texto} se traga su argumento: no produce ninguna salida en el
% PDF. El texto sigue visible en el .tex y en el .md que lo genera, y la
% lista autoritativa vive en PENDIENTES.md, seccion "Marcadores
% [VERIFICAR] en el cuerpo".
% Para que ocultarlos no equivalga a olvidarlos, cada uno deja rastro en el
% log de compilacion y al final se emite un recuento. Deliberadamente NO se
% vuelca el texto del marcador al log: varios contienen \emph y \texttt, y
% expandirlos en un \typeout rompe la compilacion.
% OJO CON LOS ESPACIOS. Estas macros NO deben terminar en \ignorespaces.
% Un marcador se escribe pegado al punto de la frase anterior
%     ...no lo evalua.\verificar{...} Si puede enunciarse...
% y \ignorespaces se comia el espacio SIGUIENTE, de modo que el PDF salia
% con "evalua.Si puede" sin separacion. Sin \ignorespaces no hay riesgo de
% espacio doble: TeX colapsa los espacios consecutivos en uno solo.
\newcounter{verificarcount}
\newcommand{\verificar}[1]{%
  \stepcounter{verificarcount}%
  \typeout{[VERIFICAR \theverificarcount] marcador oculto en pagina \thepage.}}
% Variante para los huecos de la PORTADA. Ahi un marcador invisible dejaria
% "Codigo:" seguido de nada, que se lee como error de composicion en vez de
% como dato faltante. Imprime una linea de relleno, que no afirma nada.
\newcommand{\verificarhueco}[1]{%
  \stepcounter{verificarcount}%
  \typeout{[VERIFICAR \theverificarcount] hueco de portada en pagina \thepage.}%
  \rule[-0.15em]{3.2cm}{0.4pt}}
\AtEndDocument{%
  \typeout{^^J=====================================================^^J%
    MARCADORES [VERIFICAR] EN EL CUERPO: \theverificarcount^^J%
    No se imprimen en el PDF. Detalle en PENDIENTES.md.^^J%
    =====================================================^^J}}

%--------------------- Hiperenlaces (siempre al final) ---------------
\usepackage{hyperref}
\hypersetup{
  unicode=true,
  colorlinks=true,
  linkcolor=azulesan,
  citecolor=azulesan,
  urlcolor=azulesan,
  pdftitle={Estimacion de aglomeraciones mediante vision computacional
            en el Aeropuerto Internacional Jorge Chavez},
  pdfauthor={Equipo Capstone I -- Universidad ESAN},
  pdfsubject={Capstone I -- Ingenieria de Inteligencia Artificial}
}

% Encabezado corrido: "Informe Final" a la izquierda, "Capstone Project" a
% la derecha, con filete horizontal debajo. El numero de pagina estaba antes
% en el encabezado derecho, que ahora ocupa "Capstone Project"; se baja al
% pie centrado para no perderlo (el indice remite a el).
% La portada va dentro de titlepage, que fuerza \thispagestyle{empty}: el
% encabezado aparece en todas las paginas salvo la caratula, que es lo
% habitual y lo que evita que el filete cruce el logo.
\pagestyle{fancy}
\fancyhf{}
\fancyhead[L]{\small Informe Final}
\fancyhead[R]{\small Capstone Project}
\fancyfoot[C]{\small\thepage}
\renewcommand{\headrulewidth}{0.4pt}
\renewcommand{\footrulewidth}{0pt}
% article aplica el estilo "plain" a algunas paginas sueltas; se redefine
% para que tambien lleven el encabezado y no queden sin el.
\fancypagestyle{plain}{%
  \fancyhf{}%
  \fancyhead[L]{\small Informe Final}%
  \fancyhead[R]{\small Capstone Project}%
  \fancyfoot[C]{\small\thepage}%
  \renewcommand{\headrulewidth}{0.4pt}%
  \renewcommand{\footrulewidth}{0pt}}

\onehalfspacing

%=====================================================================
\begin{document}
%=====================================================================

%--------------------------- PORTADA ---------------------------------
% LOGO DE ESAN
% No hay ningun archivo de logo en el repositorio (2026-09-13). Se deja el
% \includegraphics condicionado con \IfFileExists para que, en cuanto se
% coloque el archivo, entre solo y sin tocar este .tex.
% Formato preferible: PDF o PNG vectorial/alta resolucion, en
%   informe/imagenes/logo-esan.pdf   (o .png)
% Mientras no exista, se compone un recuadro neutro del mismo tamano para
% que la caratula no cambie de maquetacion al anadirlo.
\begin{titlepage}
\centering
% La caratula no necesita interlineado 1,5: con el logo anadido, el
% \onehalfspacing del documento la desbordaba a una segunda pagina.
\singlespacing
\vspace*{0.2cm}

\IfFileExists{imagenes/logo-esan.pdf}{%
  \includegraphics[height=2.2cm]{imagenes/logo-esan.pdf}%
}{%
  \IfFileExists{imagenes/logo-esan.png}{%
    \includegraphics[height=2.2cm]{imagenes/logo-esan.png}%
  }{%
    \framebox[5.5cm][c]{\rule{0pt}{2.2cm}\small\itshape Logo ESAN}%
  }%
}

\vspace{0.5cm}

{\scshape\large Universidad ESAN\par}
{\scshape Facultad de Ingeniería\par}
{\scshape Carrera de Ingeniería de Inteligencia Artificial\par}

\vspace{0.5cm}

{\color{azulesan}\rule{\linewidth}{1.2pt}}
\vspace{0.4cm}

{\LARGE\bfseries Estimación de aglomeraciones de personas\par}
\vspace{0.2cm}
{\LARGE\bfseries mediante visión computacional en el\par}
\vspace{0.2cm}
{\LARGE\bfseries Aeropuerto Internacional Jorge Chávez\par}

\vspace{0.4cm}
{\Large Informe Final -- Capstone Project\par}
\vspace{0.3cm}
{\large Cliente: Lima Airport Partners (LAP)\par}

\vspace{0.4cm}
{\color{azulesan}\rule{\linewidth}{1.2pt}}

\vspace{0.5cm}

{\large\bfseries Integrantes\par}
\vspace{0.4cm}

% Formato "Apellidos, Nombres (codigo)".
% Los codigos no constan en el repositorio: van con \verificarhueco, que
% imprime una linea de relleno en vez del texto del marcador.
% Solo de Moreno Ugarte constan los dos apellidos; de los otros tres el
% repositorio unicamente registra el apellido paterno.
\verificar{apellido materno de Aguilar, Rivadeneyra y Ortega: el
repositorio solo registra el apellido paterno.}%
\begin{tabular}{@{}l@{\hspace{0.6em}}l@{}}
Aguilar, Angel           & (\verificarhueco{código de Aguilar})      \\[0.2cm]
Rivadeneyra, Stephano    & (\verificarhueco{código de Rivadeneyra})  \\[0.2cm]
Ortega, José             & (\verificarhueco{código de Ortega})       \\[0.2cm]
Moreno Ugarte, Fabián    & (\verificarhueco{código de Moreno})       \\
\end{tabular}

\vfill

{\large Universidad ESAN\par}
\vspace{0.3cm}

% Ciclo academico y fecha de entrega: pendiente A-10 en PENDIENTES.md. Los dos
% valores deben salir del silabo del curso, no de una estimacion.
{\large Ciclo académico: \verificarhueco{ciclo según sílabo}\par}
\vspace{0.25cm}
{\large Fecha: \verificarhueco{fecha de entrega según sílabo}\par}
\vspace{0.45cm}

{\large Lima, Perú\par}

\end{titlepage}

%--------------------------- INDICE ----------------------------------
\pagenumbering{roman}
\tableofcontents
\clearpage

\phantomsection
\addcontentsline{toc}{section}{Nota sobre el estado de este documento}
\section*{Nota sobre el estado de este documento}

Este informe se encuentra en elaboración. Las secciones de
\emph{Antecedentes}, \emph{Marco ético y legal}, \emph{Análisis de
impacto} y el \emph{Anexo A} corresponden al frente de validación
ético-legal y documentación. Las secciones de \emph{Introducción},
\emph{Metodología}, \emph{Resultados} y \emph{Conclusiones} están a cargo
de otros integrantes del equipo y se incorporarán en la siguiente
entrega; sus apartados figuran ya con la numeración reservada.

El cuerpo del documento no contiene marcadores de relleno. Los puntos que
siguen abiertos se registran fuera de este texto, en el archivo
\texttt{PENDIENTES.md} del repositorio, agrupados según quién puede
resolverlos: decisiones del frente ético-legal, decisiones del equipo e
información que solo puede proporcionar LAP.

Dos criterios de redacción se aplicaron de forma sistemática, y conviene
enunciarlos porque explican el tono del documento:

\begin{itemize}[leftmargin=1.4em,itemsep=0.25em]
  \item \textbf{Ninguna cifra tomada de la literatura se reproduce si no
        pudo contrastarse contra el texto en prosa del artículo de
        origen.} Las cifras que solo constaban en tablas cuya conversión
        a texto se degrada fueron retiradas, aunque resultaran
        favorables al argumento.
  \item \textbf{Ninguna afirmación normativa se enuncia en indicativo si
        su respaldo documental exacto no ha sido verificado.} Donde la
        verificación está pendiente, la afirmación se formula en
        condicional o se remite expresamente a la validación del área
        legal de LAP.
\end{itemize}

El detalle de la trazabilidad de cada afirmación normativa se encuentra
en el archivo \texttt{NOTAS\_\allowbreak FUENTES.md} del repositorio, y el estado de
cumplimiento en \texttt{CHECKLIST\_\allowbreak CUMPLIMIENTO.md}.

\cajaaviso{Alcance de este documento}{Este informe es un trabajo
académico de análisis de cumplimiento. No constituye asesoría legal ni un
dictamen jurídico. Las zonas de interpretación normativa se presentan
como riesgos a validar con el área legal de LAP y, de ser necesario, con
la Autoridad Nacional de Protección de Datos Personales.}

\clearpage
\pagenumbering{arabic}

%------------------------- SECCIONES ---------------------------------
% NOTA (2026-09-11, segunda revision) -- SECCIONES 3 Y 4 REINCORPORADAS
% Los .md siguen siendo la FUENTE de las Secciones 3 y 4:
%   secciones/03_marco_normativo.md
%   secciones/04_analisis_impacto.md
% Los .tex que se incluyen aqui se GENERAN desde esos .md y no se editan
% a mano; si hay que cambiar algo de la Seccion 3 o 4, se cambia en el
% .md y se vuelve a convertir.
% Estos dos \input apuntaban antes a secciones/_archivo/, que quedo
% OBSOLETO y ya no interviene en la compilacion. Los \label de los
% archivos generados reproducen los anclajes de aquella version para que
% las referencias cruzadas del resto del informe sigan resolviendo.
% =====================================================================
% INTRODUCCION -- SECCION A CARGO DE OTRO INTEGRANTE
% Responsable: no corresponde al frente etico-legal.
% Fabian: no editar el contenido; mantener el \section y los \label para
% que las referencias cruzadas del resto del informe no se rompan.
%
% ESTADO (ver PENDIENTES.md): la decision de arquitectura (B-01/D-11) sigue
% abierta -- hay una contradiccion sin resolver entre PENDIENTES.md /
% CHECKLIST_CUMPLIMIENTO.md (la marcan abierta) y fichas_antecedentes/
% B3_ByteTrack.md (que da por acordado un alcance limitado a conteo
% agregado). Por eso esta Introduccion se redacta a proposito SIN
% comprometerse con una granularidad de salida especifica: los objetivos
% remiten esa decision a la Seccion de Metodologia y a la auditoria de
% marco_etico_legal.tex (Sec. sec:auditoria), que la resuelven cuando el
% frente tecnico la documente. Metodologia, Resultados y Conclusiones
% siguen pendientes por ese motivo y no se tocan en este cambio.
% =====================================================================
\section{Introducción}
\label{sec:introduccion}

\subsection{Contexto del problema}
\label{sec:contexto}

El Aeropuerto Internacional Jorge Chávez, operado por Lima Airport
Partners (LAP), concentra en determinadas franjas horarias y en puntos
específicos del terminal --- filtros de seguridad, migraciones,
mostradores de \emph{check-in}, salas de embarque --- un volumen de
pasajeros considerablemente mayor que el de otras zonas del recinto en el
mismo momento. Conocer dónde y cuándo se concentran las personas es
información operativa de valor para un operador aeroportuario: permite,
en principio, anticipar la apertura de mostradores adicionales, redirigir
el flujo de pasajeros o dimensionar al personal de atención con mayor
precisión que una revisión manual o periódica.

LAP cuenta ya con infraestructura de videovigilancia de circuito cerrado
en el terminal. Este proyecto explora si esa misma infraestructura puede
aprovecharse, mediante técnicas de visión computacional, para producir
una estimación automatizada de la concentración de personas por zona, sin
que ello suponga identificar, individualizar ni caracterizar
biométricamente a los pasajeros que aparecen en la imagen. Esta última
condición no es una preferencia de diseño sino un límite impuesto por la
Ley 29733 de Protección de Datos Personales y su reglamento vigente
(D.S.\ 016-2024-JUS): un dato biométrico que por sí mismo identifica al
titular constituye dato sensible, y el proyecto se plantea evitar ese
tratamiento desde el diseño en lugar de mitigarlo después. La
Sección~\ref{sec:marco} desarrolla en detalle el marco normativo
aplicable, y la Sección~\ref{sec:antecedentes} revisa la literatura
técnica que hace posible, en principio, estimar una aglomeración sin
recurrir a identificación individual.

Conviene precisar el alcance de esta afirmación desde ya, para no
sobredimensionar el problema que el proyecto atiende. El equipo no cuenta,
al momento de escribir este informe, con evidencia documentada por LAP
--- incidentes registrados de aglomeración, tiempos de espera medidos o
quejas de pasajeros --- que cuantifique el problema en el propio
terminal; esa evidencia se ha solicitado al cliente y su ausencia se deja
constancia expresa en \texttt{PENDIENTES.md}. El proyecto parte, por
tanto, de una necesidad operativa plausible y compartida por la literatura
de conteo de multitudes en espacios de alto tránsito, y no de una cifra
propia de incidentes en el Aeropuerto Jorge Chávez.

\subsection{Objetivos}
\label{sec:objetivos}

\subsubsection*{Objetivo general}

Desarrollar y evaluar un sistema de visión computacional capaz de estimar
la concentración de personas en distintas zonas del Aeropuerto
Internacional Jorge Chávez a partir de imágenes de cámaras fijas, sin
recurrir a reconocimiento facial ni a la extracción de rasgos biométricos
identificantes, y en condiciones de cumplimiento con el marco normativo
peruano de protección de datos personales.

\subsubsection*{Objetivos específicos}

\begin{enumerate}[leftmargin=1.6em,itemsep=0.35em]
  \item Revisar el estado del arte en detección, localización y conteo de
        personas en escenas congestionadas, e identificar qué enfoque
        técnico resulta más adecuado para un entorno aeroportuario y para
        las restricciones de privacidad del proyecto (desarrollado en la
        Sección~\ref{sec:antecedentes}).
  \item Determinar el marco de cumplimiento normativo aplicable al
        tratamiento de imágenes de un espacio de uso público bajo gestión
        concesionada, incluyendo el análisis de impacto y los requisitos
        de privacidad desde el diseño que ese tratamiento exige
        (desarrollado en las Secciones~\ref{sec:marco}
        y~\ref{sec:analisis}).
  \item Diseñar e implementar el componente de detección y estimación,
        con la granularidad de salida --- superficie de densidad, conteo
        agregado por zona o localización individual --- que resulte de
        aplicar el principio de proporcionalidad desarrollado en la
        Sección~\ref{sec:directiva} y la síntesis de
        antecedentes de la Sección~\ref{sec:sintesis-antecedentes}. Esta
        decisión se documentará en la Sección de Metodología cuando el
        frente técnico la cierre; hasta entonces queda registrada como
        abierta en \texttt{PENDIENTES.md} (puntos B-01 y D-11).
  \item Evaluar el desempeño del sistema mediante métricas cuantitativas y
        cualitativas, indicando en cada caso sobre qué conjunto de datos
        fue medido cada resultado --- ninguna cifra obtenida sobre
        conjuntos públicos se presentará como desempeño esperable en el
        terminal sin la validación en sitio que la
        Sección~\ref{sec:ant-p2r} señala como condición previa.
  \item Identificar los requisitos de validación adicionales que quedan
        fuera del alcance material de este Capstone --- en particular el
        acceso a datos del propio terminal y la evaluación de sesgo
        demográfico --- y dejarlos documentados como trabajo pendiente en
        lugar de darlos por resueltos (retomado en la
        Sección~\ref{sec:conclusiones}).
\end{enumerate}

\subsection{Alcance y limitaciones del proyecto}
\label{sec:alcance}

\textbf{Alcance funcional.} El proyecto cubre la estimación de la
concentración de personas por zona a partir de imágenes de cámaras fijas.
Queda expresamente fuera de su alcance cualquier forma de identificación,
reidentificación o caracterización individual de las personas detectadas:
el sistema está pensado para responder \emph{cuántas personas hay, y
dónde}, no \emph{quién es cada una}. La Sección~\ref{sec:ant-p2pnet}
señala, sin embargo, que esta distinción es una cuestión de política de
persistencia y no solo de algoritmo: el mismo modelo que localiza personas
puede, si sus salidas se retienen y se asocian entre cuadros, aproximarse
a un seguimiento de trayectorias. Por eso la granularidad de salida que
finalmente se implemente es, como se indica en el objetivo específico 3,
una decisión que debe quedar documentada por escrito antes de cerrar la
Sección de Metodología.

\textbf{Zonas cubiertas.} Las zonas concretas del terminal sobre las que
se entrena y evalúa el sistema dependen de qué cámaras y qué material
audiovisual habilite LAP para el proyecto, dato que a la fecha de este
informe no ha sido confirmado por el cliente (\texttt{PENDIENTES.md},
punto C-09). Donde este informe menciona zonas a título de ejemplo ---
filtros de seguridad, mostradores de \emph{check-in}, salas de embarque
--- lo hace como ilustración del tipo de espacio al que el sistema podría
aplicarse, no como una cobertura ya acordada con el cliente.

\textbf{Fase del proyecto y limitación de dominio.} Esta primera entrega
corresponde a Capstone I: un desarrollo y validación conceptual del
sistema sobre conjuntos de datos públicos de conteo de multitudes, sin
acceso todavía a material propio del Aeropuerto Jorge Chávez. Esto impone
la limitación más importante del proyecto en esta etapa, y conviene
enunciarla sin atenuarla porque condiciona la lectura de todo lo que
sigue: como documenta en detalle la Sección~\ref{sec:ant-p2r}, los propios
autores de la literatura revisada reconocen una brecha de dominio entre
los conjuntos públicos de entrenamiento y un entorno de despliegue nuevo,
y señalan que el desempeño se degrada de forma marcada cuando no media una
etapa de adaptación de dominio. En consecuencia, ninguna cifra de
desempeño que se reporte en la Sección de Resultados debe leerse como una
predicción de exactitud para el Aeropuerto Jorge Chávez: es, en el mejor
de los casos, una cota de referencia sobre datos públicos, sujeta a
validarse en sitio antes de comunicarse a LAP como una expectativa de
desempeño (Sección~\ref{sec:auditoria}).

\clearpage
% =====================================================================
% ANTECEDENTES -- Responsable: Fabian Moreno Ugarte
%
% Fuente: los tres PDF de ./papers/.
%
% CRITERIO APLICADO A LAS CIFRAS (no modificar sin leer NOTAS_FUENTES.md):
% solo se reproducen las cifras que constan en el TEXTO EN PROSA de cada
% articulo (resumen, introduccion o discusion de resultados). Las que
% figuraban unicamente en las tablas de resultados fueron retiradas,
% porque la conversion de esos cuadros a texto plano mezcla filas y
% columnas y una cifra mal atribuida a un conjunto de datos seria
% exactamente el tipo de afirmacion no verificable que este informe debe
% evitar. La trazabilidad esta en NOTAS_FUENTES.md, seccion "Papers".
% =====================================================================
% 14/09/2026: pasa a ser la Seccion 4 (Marco Teorico) del indice del
% modelo del curso. Contenido sin cambios; se conserva sec:antecedentes.
\section{Marco Teórico}
\label{sec:antecedentes}

El conteo de personas en escenas congestionadas es un problema con una
línea de investigación propia dentro de la visión computacional. Esta
sección revisa tres trabajos que marcan tres etapas sucesivas de esa
línea (la regresión de mapas de densidad, la detección puramente
basada en puntos y el aprendizaje semi-supervisado sobre puntos) y
discute, para cada uno, qué propone, por qué resulta pertinente para el
caso del Aeropuerto Internacional Jorge Chávez y qué limitación
reconocen sus propios autores.

Se privilegia deliberadamente la limitación \emph{declarada por los
autores} sobre la crítica externa: en un informe de cumplimiento, lo que
un proveedor de tecnología reconoce por escrito sobre su propio método
es el insumo más defendible para dimensionar el riesgo residual del
sistema.

Solo se reproducen aquí las cifras que constan en el texto en prosa de
cada artículo, es decir, en su resumen, introducción o discusión de
resultados. Las que
figuraban únicamente en las tablas de resultados se omiten: la conversión
de esos cuadros a texto plano mezcla filas y columnas, y una cifra mal
atribuida a un conjunto de datos sería precisamente el tipo de afirmación
no verificable que este informe debe evitar. Donde la comparación importa,
se citan las magnitudes relativas que los propios autores enuncian en
prosa, que es la forma en que además presentan sus contribuciones.

% ---------------------------------------------------------------------
\subsection{CSRNet: redes convolucionales dilatadas para escenas altamente congestionadas}
\label{sec:ant-csrnet}

\subsubsection*{Qué propone}

Li, Zhang y Chen~\cite{Li_2018_CVPR} proponen CSRNet, una red
completamente convolucional para reconocimiento de escenas congestionadas
compuesta por dos bloques: un \emph{front-end} que reutiliza las primeras
diez capas convolucionales de VGG-16 para extracción de características
2D, y un \emph{back-end} de convoluciones dilatadas que amplía el campo
receptivo sin recurrir a capas de \emph{pooling}. La motivación explícita
es prescindir de las arquitecturas multi-columna previas: los autores
muestran experimentalmente que las tres columnas de MCNN aprenden
características casi idénticas, de modo que la ramificación introduce
redundancia y tiempo de entrenamiento sin aportar diversidad de escala
real. La salida del modelo no es un número sino un mapa de densidad,
entrenado con pérdida euclídea contra un mapa de referencia generado
difuminando cada anotación de cabeza con un núcleo gaussiano
geométricamente adaptativo.

Los autores enuncian sus resultados en términos comparativos frente al
estado del arte del momento: sobre ShanghaiTech obtienen un 7\,\% menos de
MAE que CP-CNN en la Parte A y un 47,3\,\% menos en la Parte B. Extienden
además el método al conteo de vehículos sobre TRANCOS, donde reportan un
15,4\,\% menos de MAE que el mejor método previo.

\subsubsection*{Por qué aplica al caso LAP}

Tres propiedades del diseño son directamente relevantes para el terminal.

\begin{enumerate}[leftmargin=1.6em,itemsep=0.3em]
  \item \textbf{La salida es un mapa de densidad, no una identidad.} El
        modelo produce una superficie continua de concentración espacial
        y el conteo se obtiene integrándola. Ninguna etapa del
        \emph{pipeline} requiere extraer rasgos faciales ni construir una
        plantilla biométrica. Como se argumenta en la
        Sección~\ref{sec:auditoria}, esta característica es el principal
        insumo técnico para sostener que el tratamiento no recae sobre
        datos biométricos, aunque la conclusión depende también de qué se
        almacene aguas arriba del modelo.
  \item \textbf{La arquitectura es puramente convolucional y admite
        resoluciones variables.} Esto importa operativamente: las cámaras
        de un terminal no comparten resolución, altura ni ángulo, y el
        modelo no obliga a reescalar todo a un formato único.
  \item \textbf{La distribución espacial es informativa por sí misma.}
        Los autores ilustran el punto con tres imágenes que contienen el
        mismo número de personas pero distribuciones radicalmente
        distintas. Para LAP, saber que hay 300 personas en el hall de
        salidas es mucho menos útil que saber que 200 de ellas están
        concentradas frente a tres mostradores.
\end{enumerate}

\subsubsection*{Limitación reconocida por los autores}

La limitación más explícita aparece en la evaluación sobre UCSD, un
conjunto de escenas relativamente dispersas (entre 11 y 46 personas por
imagen): los propios autores indican que superan a la mayoría de métodos
previos \emph{excepto a MCNN en la categoría MAE}. El artículo consigna
en prosa que CSRNet alcanza 1,16 de MAE sobre ese conjunto, pero no es
esa cifra sino la excepción reconocida la que interesa aquí. Reconocen
además que la baja resolución de
los cuadros ($238\times158$ píxeles) dificulta generar un mapa de
densidad de calidad tras operaciones de \emph{pooling} sucesivas, y que
fue necesario reescalar por interpolación bilineal a $952\times632$.

Esto tiene una consecuencia directa para el proyecto: el régimen de baja
densidad no es el terreno donde el método brilla, y buena parte del día
un terminal aéreo opera precisamente en baja densidad. El método está
optimizado para el escenario que menos ocurre.

Una segunda limitación, más estructural, es que el mapa de densidad no
entrega posiciones individuales. Los autores lo asumen como parte del
paradigma; el trabajo siguiente lo convierte en su crítica central.

% ---------------------------------------------------------------------
\subsection{P2PNet: un marco puramente basado en puntos}
\label{sec:ant-p2pnet}

\subsubsection*{Qué propone}

Song et al.~\cite{Song_2021_ICCV} sostienen que localizar individuos
responde mejor a las demandas prácticas del análisis de multitudes que
simplemente contarlos, y que los métodos existentes basados en
representaciones intermedias (mapas de densidad o cajas delimitadoras
pseudo-generadas) son contraintuitivos y propensos a error. Su
propuesta, P2PNet, predice directamente un conjunto de puntos con sus
coordenadas y confianzas, sin pasar por ninguna representación
intermedia. La pieza crítica del método es la asignación de objetivos de
aprendizaje: los autores establecen que tanto el caso en que varias
propuestas se emparejan con una misma anotación como el caso inverso
confunden al modelo durante el entrenamiento y producen conteos sobre- o
sub-estimados, y resuelven el problema con un emparejamiento uno a uno
mediante el algoritmo húngaro.

El trabajo aporta además una métrica, la \emph{density Normalized Average
Precision} (nAP), que evalúa conjuntamente localización y conteo, y que
fue diseñada para no ignorar la variación de densidad entre escenas ni
dejar sin penalización las predicciones duplicadas.

En cuanto a resultados de conteo, los autores los describen en prosa de
forma comparativa: sobre ShanghaiTech Parte A reducen el MAE un 4,8\,\% y
el MSE un 12,9\,\% respecto del segundo mejor método, ADSCNet, y sobre la
Parte B logran una reducción del 2,3\,\% en MAE; sobre UCF\_CC\_50 rebajan
el MAE en 2,1 puntos frente al segundo mejor resultado, y sobre
NWPU-Crowd un 12,4\,\% frente a DM-Count. La única cifra absoluta que el
artículo enuncia fuera de sus tablas es un MAE de 85,32 sobre UCF-QNRF.
En localización, los autores señalan que el nAP$_{0.5}$ se sitúa cerca del
90\,\% en la mayoría de los conjuntos y reportan un F1/Precisión/Recall de
71,2\,\%/72,9\,\%/69,5\,\% bajo la métrica de NWPU-Crowd.

\subsubsection*{Por qué aplica al caso LAP}

La localización individual es lo que convierte una cifra en una decisión
operativa. Un conteo agregado por cámara no permite decidir dónde abrir
un mostrador adicional o hacia dónde redirigir el flujo; un conjunto de
puntos sí, porque admite agregación por zonas definidas por LAP (colas de
\emph{check-in}, filtro de seguridad, migraciones) sin reentrenar el
modelo.

Ahora bien, esta misma capacidad es la que eleva el perfil de riesgo del
sistema y obliga a un análisis más fino en la Sección~\ref{sec:marco}. Un
punto con coordenadas no es un dato biométrico, pero sí es un dato
espacio-temporal referido a una persona física presente en un lugar y
momento determinados. Si esos puntos se persisten y se encadenan en el
tiempo, el sistema se acerca a un seguimiento de trayectorias, que es una
finalidad distinta de la estimación de aglomeraciones y que requeriría su
propia base de legitimación.

Queda así un riesgo a validar con el cliente. La diferencia entre
\emph{estimar una aglomeración} y \emph{seguir a una persona} no es una
diferencia de algoritmo sino de política de persistencia. Con la misma
salida de P2PNet, retener los puntos y asociarlos entre cuadros produce
trayectorias individuales. Este límite debe fijarse por diseño y por
contrato, no dejarse a la configuración: una vez que el sistema está en
producción, permitir la persistencia es un cambio de parámetro, mientras
que prohibirla requiere volver sobre el diseño.

\subsubsection*{Limitación reconocida por los autores}

Los autores reconocen tres limitaciones de manera explícita.

Primero, en UCF-QNRF admiten que su precisión \enquote{no es tan
competitiva} frente a ADSCNet, aunque sea superior en todos los demás
conjuntos. Segundo, en NWPU-Crowd atribuyen su desempeño inferior en la
métrica MAE[S] a que sus predicciones se basan en un único mapa de
características de una sola escala, elegido por simplicidad; señalan que
incorporar fusión multiescala sería una mejora natural. Tercero, y más
relevante para el caso LAP, observan que la precisión es notablemente
menor bajo el umbral más estricto nAP$_{0.05}$, y explican que a esa
exigencia de localización los efectos de las desviaciones de etiquetado
empiezan a hacerse aparentes. El artículo no cuantifica esa caída fuera
de sus tablas; el contraste que sí enuncia en prosa es que el
nAP$_{0.5}$ ronda el 90\,\% mientras que ya bajo el umbral intermedio
nAP$_{0.25}$ la precisión se mantiene por encima del 55\,\%.

Esta última observación es un límite \emph{del suelo de verdad}, no del
modelo: por debajo de cierta escala, la propia anotación humana es
ruidosa. Cualquier compromiso de exactitud que el equipo ofrezca a LAP
debe formularse a nivel de zona agregada y no a nivel de posición
individual, porque a nivel individual la referencia contra la cual se
mediría no es confiable.

% ---------------------------------------------------------------------
\subsection{P2R: pérdida punto-a-región para conteo semi-supervisado}
\label{sec:ant-p2r}

\subsubsection*{Qué propone}

Lin, Zhao y Chan~\cite{Lin_2025_CVPR} parten de una constatación
económica: entrenar un contador basado en puntos exige anotar cientos o
miles de puntos por imagen, y ese costo ha frenado la adopción del
enfoque. Integran entonces los métodos basados en puntos en un marco
semi-supervisado de pseudo-etiquetado, en el que un modelo maestro genera
etiquetas para imágenes no anotadas y un modelo estudiante aprende de
ellas.

Al implementarlo detectan un problema específico: la confianza asociada a
una pseudo-etiqueta no logra propagarse a los píxeles de fondo bajo el
esquema punto-a-punto. Para diagnosticarlo diseñan un mapa de activación
específico por punto (PSAM), y observan que los píxeles vecinos de cada
píxel de primer plano quedan sobre-activados, de modo que el
decodificador los interpreta como instancias distintas. La solución
propuesta es sustituir el emparejamiento punto-a-punto por uno
punto-a-región (P2R): en lugar de detectar un punto por peatón, se
segmenta una región local, de forma que todos los píxeles de esa región
comparten la confianza del pseudo-punto correspondiente. Como efecto
colateral, el algoritmo húngaro deja de ser necesario, lo que reduce el
costo de cómputo de la pérdida.

Los experimentos se realizan sobre cuatro conjuntos (ShanghaiTech A y
B, UCF-QNRF y JHU-Crowd++) bajo protocolos de 5\,\%, 10\,\% y 40\,\%
de datos etiquetados. La afirmación central que los autores enuncian en
prosa es que entrenar con el 5\,\% de los datos usando P2R supera al
aprendizaje totalmente supervisado con el 10\,\%, y resulta equivalente a
otros métodos semi-supervisados que emplean ese mismo 10\,\%. Como
referencia del punto de partida, el artículo consigna que el modelo
entrenado únicamente con ese 5\,\% de datos etiquetados alcanza un MAE de
93,7 y un MSE de 155,2 sobre ShanghaiTech Parte A. Bajo supervisión
completa (100\,\% de etiquetas) P2R supera a P2PNet en los cuatro
conjuntos.

La supresión del algoritmo húngaro tiene además un efecto que los autores
miden y describen en prosa: sobre una imagen de $576\times960$ con 775
puntos anotados, el cálculo de la pérdida punto-a-punto requiere 0,4307
segundos frente a 0,0064 segundos de P2R, cerca de 68 veces más rápido.
No es un detalle menor para el proyecto, porque el costo de cómputo del
entrenamiento tiene una lectura ambiental que se retoma en la
Sección~\ref{sec:impacto-ambiental}.

\subsubsection*{Por qué aplica al caso LAP}

Este es, de los tres, el trabajo con la implicancia legal más favorable,
y conviene ser explícito sobre por qué.

El principio de minimización exige tratar solo los datos que resulten
necesarios para la finalidad. Un método que alcanza desempeño competitivo
con el 5\,\% de los datos etiquetados permite reducir en proporción
similar el volumen de imágenes del terminal que deben ser revisadas y
anotadas por personas. Menos imágenes anotadas significa menos exposición
de imágenes de pasajeros reales a anotadores humanos, menos copias del
material fuera del entorno controlado y una superficie de brecha
considerablemente menor. La eficiencia de etiquetado deja de ser una
ventaja de costo y se convierte en un argumento de cumplimiento.

Además, los experimentos de adaptación de dominio no supervisada del
artículo abordan de frente el problema que el proyecto tendrá que
resolver: ningún conjunto público de conteo de multitudes contiene
escenas de un terminal aéreo peruano.

\subsubsection*{Limitación reconocida por los autores}

Los autores documentan un caso de fallo atribuido al desenfoque producido
por el movimiento de la multitud, y lo presentan como tal dentro de su
comparación visual de predicciones. El método tiene, por tanto, un modo
de fallo reconocido y ligado a una condición, la multitud en movimiento,
que en un terminal aéreo es la norma y no la excepción. El artículo
no cuantifica en prosa la magnitud del error en ese caso concreto, de
modo que aquí se recoge la existencia del modo de fallo y no su tamaño.

Más importante aún, los propios autores registran la brecha de dominio.
Al transferir el modelo entre conjuntos sin adaptación de dominio el
desempeño se degrada de forma marcada, y son ellos quienes señalan
explícitamente que ese pobre desempeño subraya la brecha entre los datos
de origen y los de destino; cuando se incorporan al entrenamiento datos
no etiquetados del dominio de destino, los errores de estimación se
reducen de forma significativa. El artículo enuncia esa diferencia en
prosa sin consignarla numéricamente fuera de sus tablas.

Que los propios autores
adviertan una brecha de dominio de esta naturaleza implica que las
métricas publicadas sobre conjuntos públicos \textbf{no son
extrapolables} al Aeropuerto Jorge Chávez. Presentar a LAP una cifra de
exactitud tomada de la literatura, sin validación en el terminal, sería
una afirmación sin respaldo. Este informe recomienda, en consecuencia, que ninguna promesa cuantitativa de desempeño figure en un
entregable al cliente antes de una validación en sitio, aun cuando eso
implique presentar el proyecto sin cifras atractivas.

% ---------------------------------------------------------------------
\subsection{Síntesis: cómo se articulan los tres trabajos}
\label{sec:sintesis-antecedentes}

Los tres artículos no son alternativas entre las cuales elegir: forman
una progresión en la que cada uno resuelve una limitación del anterior y,
al hacerlo, desplaza el problema hacia el siguiente.

CSRNet~\cite{Li_2018_CVPR} resuelve el problema de \emph{cuántos} con una
arquitectura simple y entrenable de extremo a extremo, pero entrega una
superficie de densidad de la que no se puede extraer dónde está cada
persona. P2PNet~\cite{Song_2021_ICCV} resuelve el problema de
\emph{dónde} eliminando la representación intermedia y prediciendo puntos
directamente, pero paga ese avance con un costo de anotación que crece
con la densidad de la escena. P2R~\cite{Lin_2025_CVPR} resuelve el
problema del \emph{costo de anotación} llevando el enfoque por puntos al
régimen semi-supervisado, y de paso elimina el algoritmo húngaro que
P2PNet había introducido.

Vista desde el análisis de cumplimiento, la progresión tiene una lectura
distinta y en parte inversa. Cada paso hacia adelante en capacidad
técnica es también un paso hacia un perfil de riesgo distinto:

% 14/09/2026: longtable en lugar de tabularx dentro de center. Sin salto de
% pagina antes del Marco Teorico la tabla no cabia al pie de la pagina y
% dejaba un tercio en blanco; longtable puede partirse entre filas y repite
% la cabecera. Mismas columnas y mismo texto.
{\small
\setlength{\LTleft}{0pt}\setlength{\LTright}{0pt}
\begin{longtable}{@{}L{2.5cm}L{3.0cm}L{\dimexpr\linewidth-5.5cm-4\tabcolsep\relax}@{}}
\toprule
\textbf{Trabajo} & \textbf{Avance técnico} & \textbf{Desplazamiento del riesgo} \\
\midrule
\endfirsthead
\toprule
\textbf{Trabajo} & \textbf{Avance técnico} & \textbf{Desplazamiento del riesgo} \\
\midrule
\endhead
\bottomrule
\endlastfoot
CSRNet \cite{Li_2018_CVPR} &
Mapa de densidad sin representación individual &
Perfil de riesgo más bajo: la salida es agregada por construcción. El
riesgo se concentra aguas arriba, en qué se hace con el video original. \\
\addlinespace
P2PNet \cite{Song_2021_ICCV} &
Localización individual por puntos &
El riesgo se traslada a la salida del modelo: coordenadas por persona,
persistibles y encadenables en trayectorias. Exige control de
persistencia. \\
\addlinespace
P2R \cite{Lin_2025_CVPR} &
Desempeño competitivo con 5\,\% de etiquetas &
Reduce el riesgo del proceso de anotación (menos imágenes reales
expuestas a personas), pero introduce dependencia de pseudo-etiquetas
cuya calidad no es auditable caso por caso. \\
\end{longtable}
}

De esta lectura conjunta se desprenden tres criterios que el análisis de cumplimiento traslada al desarrollo
del sistema:

\begin{enumerate}[leftmargin=1.6em,itemsep=0.35em]
  \item \textbf{La granularidad de la salida es una decisión de
        cumplimiento, no solo de producto.} Si la finalidad declarada a
        LAP es estimar aglomeraciones, una salida agregada por zona
        satisface esa finalidad; una salida de coordenadas individuales
        la excede. El principio de proporcionalidad sugiere preferir la
        mínima granularidad que resuelva el problema, aunque el modelo
        sea capaz de más.
  \item \textbf{La eficiencia de etiquetado de P2R es un argumento de
        minimización aprovechable.} Conviene documentarla como tal en el
        análisis de riesgo y no solo como una ventaja de costo.
  \item \textbf{Ninguna cifra de la literatura es transferible al
        terminal sin validación local.} La propia evidencia de
        \cite{Lin_2025_CVPR} sobre brecha de dominio lo demuestra.
\end{enumerate}

Esta síntesis queda deliberadamente abierta en un punto. La elección de arquitectura del sistema no consta por escrito al cerrar esta
versión del informe, y por esa razón la auditoría de la
Sección~\ref{sec:auditoria} se formula de manera condicional sobre los
tres enfoques en lugar de sobre uno. Cuando esa elección quede
documentada, corresponde cerrar esta síntesis indicando cuál de los tres
se adopta y con qué granularidad de salida, porque de la granularidad
(mapa de densidad, puntos individuales o agregación por zona)
depende buena parte del examen de proporcionalidad de la
Sección~\ref{sec:proporcionalidad}.

\clearpage
%=====================================================================
%  SECCION 3 -- MARCO NORMATIVO APLICABLE
%  Responsable: Fabian Moreno Ugarte
%
%  GENERADO desde secciones/03_marco_normativo.md, que sigue siendo la
%  FUENTE de esta seccion. No editar aqui: editar el .md y reconvertir.
%
%  Version comprimida (2026-09-14): el desarrollo completo esta en el
%  Anexo B, generado desde secciones/anexo_b_marco_normativo.md.
%  Los marcadores label sostienen las referencias cruzadas del resto
%  del informe. La numeracion 3.1-3.4 es vinculante.
%=====================================================================

\section{Marco normativo aplicable}
\label{sec:marco}

Esta sección identifica el régimen jurídico que rige el tratamiento de
imágenes de personas mediante videovigilancia y visión computacional en
un terminal aeroportuario, y traduce cada obligación en un requisito
verificable sobre el sistema propuesto. Las zonas en que la norma admite
más de una lectura razonable se presentan como riesgos a validar y no
como conclusiones: el análisis lo elabora un equipo de estudiantes de
ingeniería y no un despacho legal, y la calificación jurídica definitiva
corresponde al área legal de Lima Airport Partners y, en su caso, a la
Autoridad Nacional de Protección de Datos Personales. Aquí se recoge lo
esencial; el desarrollo completo, con el articulado citado literalmente,
figura en el Anexo B.

% ---- 3.1 ----
\subsection{Régimen normativo aplicable}
\label{sec:directiva}
\label{sec:dl1218}

El punto de partida es el artículo 2 de la Constitución, que en sus
incisos 6 y 7 protege la intimidad personal y familiar y la imagen
propia. Captar la imagen de una persona en un espacio abierto al público
no es por ello ilícito, pero constituye una injerencia en un derecho
fundamental y debe superar un examen de proporcionalidad.

La norma central es la Ley 29733, Ley de Protección de Datos Personales.
Su reglamento vigente es el Decreto Supremo 016-2024-JUS, publicado en
\emph{El Peruano} el 30 de noviembre de 2024, que derogó el D.S.
003-2013-JUS. Su entrada en vigor se habría producido el 30 de marzo de
2025, dato que se enuncia con reserva: es el único de los relativos a
este decreto que no se ha contrastado todavía contra la fuente oficial,
y es el que fija el régimen transitorio, de modo que debe confirmarse
antes de apoyarse en él. En cualquier caso, los procedimientos y
cláusulas que LAP tenga construidos bajo el reglamento anterior deben
revisarse. \verificar{únicamente la fecha de entrada en vigor, el 30 de
marzo de 2025, que es el dato que fija el régimen transitorio. La
expedición, la publicación y la disposición derogatoria sí se
contrastaron contra la separata de \emph{El Peruano} del 30 de noviembre
de 2024.}

Para videovigilancia rige además la Directiva 01-2020-JUS/DGTAIPD, que
sigue vigente: la disposición derogatoria del D.S. 016-2024-JUS deroga
únicamente el D.S. 003-2013-JUS. A ese cuerpo se superpone el régimen de
seguridad ciudadana del Decreto Legislativo 1218 y de la Ley 30120,
reglamentados por el Decreto Supremo 007-2020-IN. La denominación
oficial completa y la fecha de publicación del D.L. 1218 y de la Ley
30120 no se han contrastado todavía contra su texto publicado, a
diferencia del D.S. 007-2020-IN, por lo que ambas normas se citan por su
denominación usual y esos dos datos deben confirmarse antes de
reproducirse en un documento dirigido al cliente. El artículo 3 de ese
reglamento comprende los establecimientos comerciales abiertos al
público con un aforo de cincuenta personas o más, umbral que un terminal
aeroportuario supera con holgura: su aplicación al Aeropuerto Jorge
Chávez es el supuesto de partida y no una hipótesis remota.
\verificar{denominación oficial completa y fecha de publicación del D.L.
1218 y de la Ley 30120. El D.S. 007-2020-IN sí se contrastó contra el
texto publicado en \emph{El Peruano}.}

La zona gris principal es la calificación del propio terminal, un
espacio concesionado a un operador privado, abierto al público y con
funciones de interés público, que no encaja limpiamente ni como dominio
público ni como recinto privado. De esa calificación depende qué régimen
es exigible y con qué intensidad, y es el primer punto a elevar al área
legal del cliente.

% ---- 3.2 ----
\subsection{Datos tratados, legitimación y proporcionalidad}
\label{sec:proporcionalidad}

El sistema trata cuatro tipos de dato con regímenes distintos. El
fotograma contiene rostros identificables y es dato personal en el
sentido del artículo 2 numeral 4 de la Ley 29733. El punto con
coordenadas que producen los métodos basados en puntos
\cite{Song_2021_ICCV,Lin_2025_CVPR}, desligado del fotograma, no
identifica por sí solo a nadie. Asociado a una marca de tiempo y
encadenado con otros del mismo individuo, en cambio, reconstruye una
trayectoria, y esa diferencia depende de la política de persistencia y
no del algoritmo. El dato biométrico, por último, es sensible conforme
al numeral 5 solo cuando por sí mismo permite identificar al titular:
una plantilla facial lo es; un mapa de densidad o un conjunto de
coordenadas de cabezas, no. Que la salida del modelo no sea biométrica
no significa que no lo sea la entrada, de modo que almacenar o
transmitir fotogramas cambia el análisis aunque el modelo nunca ejecute
reconocimiento facial.

La presentación preliminar del proyecto habla de «personas
anonimizadas», y la documentación del prototipo califica su
identificador por persona de «no reversible». Un identificador
persistente que permite reagrupar las apariciones de una persona no es
anonimización, cuyo procedimiento es irreversible (numeral 12), sino
disociación, cuyo procedimiento es reversible (numeral 13). La
diferencia no es cosmética: el dato disociado sigue sujeto al régimen
general de la ley.

El consentimiento «previo, informado, expreso e inequívoco» del artículo
13.5 es inalcanzable en un terminal, porque no puede negarse sin
renunciar al servicio. La base habrá de buscarse entre los supuestos del
artículo 14, que la Directiva recoge en su numeral 6.3. El que la
práctica invoca con más frecuencia alude a los intereses legítimos del
titular de los datos, no a los del responsable como en el régimen
europeo, y de esa diferencia depende que ampare una finalidad operativa.
Qué base invoca LAP es información que solo el cliente puede aportar. La
finalidad, en cambio, no deja margen: los artículos 6 y 28 numeral 4
impiden extender el tratamiento a otra finalidad distinta de la
declarada. Medir la permanencia de personas en zonas comerciales, que el
propio proyecto llegó a contemplar como aplicación posible, no es una
variante de estimar aglomeraciones sino una finalidad distinta, con su
propia base de legitimación y su propio deber de información.

El artículo 7 exige que el tratamiento sea «adecuado, relevante y no
excesivo» para su finalidad. La estimación automática de densidad es
idónea, y si la finalidad es estimar aglomeraciones, el medio menos
invasivo que la alcanza produce conteos agregados sin identificar a
nadie. La ponderación en sentido estricto requiere, sin embargo,
evidencia sobre la magnitud real del problema que solo LAP puede
aportar, y mientras falte el examen descansa sobre una premisa no
acreditada. De este principio depende la decisión de mayor consecuencia
normativa: a igualdad de utilidad operativa, el mapa de densidad de
CSRNet \cite{Li_2018_CVPR} o la agregación por zona superan el examen de
necesidad mejor que la localización individual por puntos, cuya licitud
depende de que las coordenadas no se persistan ni se encadenen. La
elección del equipo no consta todavía por escrito.

% ---- 3.3 ----
\subsection{Privacidad desde el diseño y requisitos del sistema}
\label{sec:privacidad-diseno}

El D.S. 016-2024-JUS recoge en el Artículo IX de su Título Preliminar el
principio de responsabilidad proactiva: el responsable debe estar en
condiciones de acreditar que cumple, adoptando las medidas antes de
iniciar el tratamiento. De ahí se sigue la privacidad por diseño y por
defecto, que en este régimen es exigencia y no buena práctica, aunque la
norma no fija una lista mínima de medidas. El proyecto la traduce en
diez requisitos: procesamiento en el borde (M-01); no persistencia del
fotograma (M-02); salida agregada por zona, sin coordenadas individuales
persistidas (M-03); difuminado irreversible de rostros si algún
fotograma debe conservarse (M-04); plazo de retención con purga
automática (M-05); control de acceso por roles con auditoría (M-06);
cifrado en tránsito y en reposo (M-07); carteles informativos conforme
al Anexo A (M-08); linaje de datos documentado (M-09); y supervisión
humana, sin acciones automáticas sobre personas (M-10). Son requisitos y
no descripción del sistema actual: el prototipo respeta M-02, pero
\textbf{no cumple M-03}, como se expone en 3.4.

Sobre la retención, el artículo 8 de la Ley 29733 ordena conservar los
datos solo por el tiempo necesario, y la Directiva fija para las
imágenes un máximo de treinta a sesenta días. El artículo 17.2 del D.S.
007-2020-IN obliga, en cambio, a conservar las grabaciones un mínimo de
cuarenta y cinco días. Ambos regímenes solo son compatibles porque
recaen sobre finalidades distintas, y de ahí la consecuencia de diseño:
el sistema de conteo debe ser un subsistema separado del CCTV de
seguridad, con su propia finalidad, base y plazo, para no heredar una
obligación de conservación que no necesita.

Los artículos 18 a 25 reconocen los derechos de información, acceso,
rectificación, supresión, oposición, tutela e indemnización. El artículo
23, que protege frente a decisiones significativas basadas únicamente en
un tratamiento automatizado, es el fundamento directo de M-10: mientras
el sistema emita alertas agregadas que interpreta un operador, no se
activa. El deber de información se cumple mediante el cartel del numeral
6.11 de la Directiva, que junto con la proporcionalidad sostiene la
licitud del tratamiento; sus campos de identidad del responsable, base
de legitimación, canal de derechos y plazo de conservación solo puede
completarlos LAP.

Cualquier proveedor que trate datos por cuenta de LAP es encargado
sujeto al artículo 30. Si el procesamiento se ejecuta fuera del Perú,
hay flujo transfronterizo sujeto al artículo 15, y el Artículo VI del
Título Preliminar del reglamento extiende el régimen peruano a ciertos
tratamientos realizados en el exterior. Por eso la elección entre
inferencia \emph{on-premise} y en nube externa es, a juicio de este
informe, la de mayor impacto legal del proyecto; M-01 la anticipa en el
sentido más protector, pero la decisión no consta por escrito.

% ---- 3.4 ----
\subsection{Auditoría del prototipo: ¿el sistema trata datos biométricos?}
\label{sec:auditoria}
\label{sec:licencias}

La pregunta que decide el régimen aplicable es si en algún punto se
extrae de la imagen una característica que por sí sola permita
identificar a una persona. Bajo CSRNet \cite{Li_2018_CVPR} no se
construye ninguna plantilla y la salida es agregada por construcción.
Bajo P2PNet o P2R \cite{Song_2021_ICCV,Lin_2025_CVPR} tampoco hay
plantillas, pero la conformidad queda condicionada a que las coordenadas
no se persistan. Bajo el seguimiento multicámara con reidentificación
por apariencia la calificación cambia: tanto si el vector de apariencia
es dato sensible como si no lo es, el sistema reconstruye trayectorias
individuales y excede lo que la finalidad requiere conforme al artículo
7.

Contrastado con el prototipo que el equipo ha desarrollado, el resultado
es mixto. \textbf{La reidentificación por apariencia no se confirma.} La
fusión entre cámaras se resuelve por asociación espacio-temporal sobre
un plano común, y solo ante ambigüedad interviene un descriptor de
desempate (histograma de color del torso y proporción de la caja) que no
es biométrico y no se persiste. La decisión va en la dirección que exige
uno de los requisitos no funcionales que el proyecto se fijó, es decir,
de las condiciones que el sistema debe cumplir con independencia de su
funcionalidad: el que prohíbe derivar de la imagen rasgos que permitan
identificar a una persona, ni siquiera de forma indirecta (identificado
en el proyecto como RNF-01). Su texto literal y el documento en que
consta no figuran en la documentación disponible al cerrar esta versión,
de modo que se recoge según lo reportado por el equipo y esa conformidad
queda sujeta a contrastarlo con su fuente. \verificar{texto literal del
RNF-01 y documento en que consta. No figura en el repositorio al cerrar
esta versión, de modo que su formulación se recoge según lo reportado
por este frente y no según su fuente.} \textbf{Lo que se persiste, en
cambio, no es conforme.} El esquema de base de datos guarda, por cada
detección, un identificador de persona, la cámara, las coordenadas, el
instante y la zona: trayectorias individuales escritas en disco, que
incumplen M-03.

Del contraste resultan cuatro hallazgos no conformes. El primero es la
ausencia de la Evaluación de Impacto en Protección de Datos, exigible
antes de iniciar tratamientos de alto riesgo como la videovigilancia
masiva; es el más grave en lo formal y el más fácil de subsanar, porque
su contenido sustantivo está en la Sección 4. El segundo es el
incumplimiento de M-03, un defecto de implementación que el propio
equipo puede corregir agregando por zona antes de escribir o fijando una
retención muy corta sobre la tabla de posiciones. El tercero y el cuarto
son ausencias de evidencia que requieren acceso a material del terminal:
no hay validación de exactitud en el dominio de destino, y ninguno de
los trabajos revisados evalúa el sesgo demográfico de forma desagregada.
Quedan además pendientes la política de retención, la localización del
procesamiento y la trazabilidad de licencias de los conjuntos de datos.

\clearpage
%=====================================================================
%  SECCION 4 -- ANALISIS DE IMPACTO Y PRIVACIDAD DESDE EL DISENO
%  Responsable: Fabian Moreno Ugarte
%
%  GENERADO desde secciones/04_analisis_impacto.md, que sigue siendo la
%  FUENTE de esta seccion. No editar aqui: editar el .md y reconvertir.
%
%  Los marcadores label reproducen los anclajes de la version archivada. La
%  numeracion 4.1-4.5 es vinculante.
%=====================================================================

\section{Análisis de impacto y privacidad desde el diseño}
\label{sec:analisis}

% ---- 4.1 ----
\subsection{Metodología del análisis de impacto}
\label{sec:eipd}

Esta sección evalúa el impacto del sistema propuesto sobre los derechos
de las personas observadas y sobre el entorno en que se despliega. Su
estructura no es arbitraria: reproduce deliberadamente la de una
Evaluación de Impacto en Protección de Datos, esto es, identificar el
tratamiento, valorar su necesidad y proporcionalidad, identificar los
riesgos para los derechos de los titulares y proponer medidas de
mitigación.

Esa correspondencia no es una analogía metodológica. Como se expuso en
la Sección 3.4, el D.S. 016-2024-JUS exige una evaluación de impacto
antes de iniciar tratamientos de alto riesgo, y la videovigilancia
masiva figura expresamente entre los supuestos que la activan. El
sistema parte de la captación masiva y continua de imágenes de personas
en las zonas comunes de un terminal aeroportuario, de modo que el
supuesto se activa con independencia de que el modelo produzca conteos
agregados y no identifique a nadie. Esta sección es, por tanto, el
insumo sustantivo de un documento que el proyecto está obligado a
producir, y no un ejercicio analítico de acompañamiento.

Conviene anticipar una tentación previsible y descartarla expresamente:
argumentar que, como la salida es agregada y no hay reconocimiento
facial, el sistema no alcanza el umbral de alto riesgo y la evaluación
no sería exigible. Este informe recomienda no construir el argumento en
esa dirección, por dos razones. La primera es que el supuesto que se
activa es el de captación masiva, y esa captación existe con
independencia de lo que el modelo haga después. La segunda es que la
evaluación de impacto es precisamente el documento en el que las
decisiones de diseño de este informe (minimización, no persistencia,
salida agregada, procesamiento local) se registran como medidas de
mitigación y quedan acreditadas ante la autoridad. Renunciar a ella
sería renunciar al lugar donde el trabajo de cumplimiento se vuelve
exigible y demostrable.

Un criterio gobierna la redacción de todo lo que sigue: se prefiere
declarar un impacto como no cuantificado antes que estimarlo sin base.
Varias magnitudes que naturalmente irían aquí (número de pasajeros
afectados, costo evitado, consumo energético, tasa de error por grupo
demográfico) dependen de información operativa que solo LAP posee o de
mediciones que el equipo aún no ha realizado. Donde ocurre, se indica
expresamente en lugar de rellenarse con una cifra plausible.

% ---- 4.2 ----
\subsection{Sesgos del sistema}

El sistema puede fallar de manera desigual, y conviene distinguir tres
fuentes de desigualdad que operan por mecanismos distintos y admiten
mitigaciones distintas. En los tres casos se indica qué reconocen los
propios autores de la literatura revisada, porque la solidez del
análisis depende de no atribuirles limitaciones que no declararon ni de
silenciar las que sí.

\textbf{Brecha de dominio.} Es la limitación mejor documentada y la más
grave para esta fase del proyecto. Los autores de P2R
\cite{Lin_2025_CVPR} reconocen expresamente que el desempeño de los
métodos de conteo se degrada de forma marcada cuando el entorno de
despliegue difiere del de entrenamiento, y que sin una etapa de
adaptación de dominio las cifras publicadas no se trasladan. Documentan
además un modo de fallo concreto del paradigma punto a punto: el
vecindario de cada píxel de primer plano queda sobre-activado y se lee
como instancia distinta, lo que produce sobreestimación severa del
conteo. La consecuencia para el proyecto es directa y condiciona la
lectura de toda la Sección de Resultados: esta primera entrega se
desarrolla sobre conjuntos públicos, sin acceso todavía a material del
Aeropuerto Jorge Chávez, de modo que ninguna cifra que se reporte debe
leerse como predicción de exactitud para el terminal. Es, en el mejor de
los casos, una cota de referencia sobre datos públicos, sujeta a
validación en sitio antes de comunicarse a LAP como expectativa de
desempeño.

\textbf{Sesgo demográfico.} Aquí la situación es distinta y debe
enunciarse con precisión: no es que los autores reconozcan un sesgo
demográfico, sino que ninguno de los trabajos revisados reporta
desempeño desagregado por características demográficas de las personas
contadas. No se trata de una omisión de este informe sino de una
ausencia en la literatura consultada, y tiene consecuencias concretas en
un aeropuerto internacional, donde la población observada es
considerablemente más heterogénea que la de los conjuntos de
entrenamiento disponibles. Los factores que razonablemente podrían
inducir desempeño desigual son identificables aunque no estén
cuantificados. La estatura es el primero: un contador entrenado sobre
anotaciones de cabezas puede subcontar sistemáticamente a niños, sobre
todo en escenas densas donde quedan ocluidos por adultos. La vestimenta
y los cubrimientos de cabeza (sombreros, gorras, capuchas, prendas
religiosas) alteran la apariencia de la región de la cabeza, que es
precisamente la señal que estos métodos explotan. El equipaje voluminoso
y las sillas de ruedas presentan siluetas distintas de las típicas en
los conjuntos de origen. Y la composición demográfica de esos conjuntos,
procedentes mayoritariamente de escenas de otras regiones del mundo, no
ofrece garantía alguna de neutralidad respecto de la población
retratada: la brecha documentada en \cite{Lin_2025_CVPR} es geográfica y
de escena, y no hay evidencia de que sea neutral respecto de quién
aparece en la imagen.

\textbf{Sesgo de anotación.} Es el menos visible de los tres y el que
los autores sí documentan de forma explícita. Los responsables de P2PNet
\cite{Song_2021_ICCV} señalan que su métrica de localización más
exigente, el nAP con umbral 0,05, cae de forma notable por efecto de las
desviaciones de etiquetado, es decir, por la imprecisión con que los
anotadores humanos sitúan el punto sobre la cabeza. El sesgo no procede
aquí del modelo sino de la verdad de referencia contra la que se lo
entrena y se lo mide. En la misma línea, los autores de CSRNet
\cite{Li_2018_CVPR} construyen su mapa de referencia difuminando cada
anotación con un núcleo gaussiano geométricamente adaptativo, lo que
incorpora un supuesto sobre la relación entre distancia a la cámara y
tamaño aparente de la cabeza; cuando ese supuesto no se cumple
(perspectivas oblicuas, techos altos, cámaras de gran angular, todas
ellas habituales en un terminal), la referencia misma queda sesgada. Los
propios autores reconocen además que sobre UCSD, un conjunto de escenas
relativamente dispersas, superan a la mayoría de métodos previos excepto
a MCNN en MAE, y que la baja resolución de esos cuadros les obligó a
reescalar por interpolación bilineal. Esto último importa operativamente
más de lo que parece: el régimen de baja densidad no es el terreno donde
el método brilla, y buena parte del día un terminal aéreo opera
precisamente en baja densidad.

Conviene ser preciso sobre la gravedad de estos sesgos, para no
exagerarla ni minimizarla. El sistema no toma decisiones sobre personas
individuales: no deniega embarques ni señala sospechosos, de modo que un
error de conteo no perjudica directamente a la persona mal contada. Sin
embargo, si el subconteo se concentra en un grupo, las zonas del
terminal donde ese grupo se concentra reciben sistemáticamente menos
recursos: menos personal y menos apertura de mostradores. El daño no es
individual sino distributivo, y es real aunque sea indirecto. De ahí se
sigue una recomendación concreta: que la campaña de validación en sitio
incluya una evaluación de desempeño desagregada, al menos por franja de
estatura y por presencia de equipaje voluminoso. La recomendación es
condicional en sentido estricto, porque sin acceso a material del
terminal no es realizable, y esa dependencia refuerza la solicitud de
acceso planteada en la Sección 3.4.

% ---- 4.3 ----
\subsection{Sesgo de automatización y factor humano}
\label{sec:impacto-social}

Un riesgo distinto de los anteriores, y menos discutido, es la
sustitución progresiva del criterio humano por el del sistema. Si el
personal de operaciones llega a confiar en el tablero más que en su
propia observación, un fallo del sistema en un momento crítico produce
un efecto peor que su ausencia, porque la atención humana ya se retiró.
El fenómeno no depende de que el modelo funcione mal: depende de que
funcione bien durante el tiempo suficiente para que nadie lo cuestione.

Este riesgo tiene un anclaje normativo directo. Como se expuso en la
Sección 3.3, el artículo 23 de la Ley 29733 reconoce el derecho al
tratamiento objetivo, esto es, a no verse sometido a una decisión con
efectos jurídicos o significativos sustentada únicamente en un
tratamiento automatizado. La medida M-10, conforme a la cual el sistema
emite alertas y no ejecuta acciones automáticas sobre personas, responde
a esa exigencia y mantiene al sistema fuera del supuesto del artículo
23. Pero M-10 es insuficiente por sí sola, porque describe la
arquitectura y no la práctica: un operador que aprueba mecánicamente
cuanto el tablero propone reintroduce de hecho la decisión automatizada
que la medida pretendía evitar, aunque formalmente medie una acción
humana.

De ahí que la capacitación del personal deba tratarse como medida de
mitigación y no como accesorio de implantación. Los operadores necesitan
conocer las limitaciones del modelo, y en particular tres de las
documentadas en esta sección: que el desempeño sobre datos públicos no
predice el desempeño en el terminal mientras no exista validación en
sitio; que existe un modo de fallo por sobreestimación severa
documentado en \cite{Lin_2025_CVPR}; y que el régimen de baja densidad,
habitual durante buena parte de la jornada, no es aquel en que estos
métodos se comportan mejor. Un operador que conoce estos tres puntos
interpreta una alerta anómala como lo que probablemente es; uno que no
los conoce, la ejecuta.

Existe además un impacto sobre el comportamiento de las personas
observadas que conviene dejar planteado con la cautela que corresponde.
Se sostiene habitualmente que saberse observado modifica la conducta;
este informe no puede afirmarlo como hecho acreditado, porque las
referencias validadas del proyecto son los artículos técnicos de
antecedentes y no se ha incorporado literatura sobre efectos
conductuales de la videovigilancia. Formulado en condicional: si ese
efecto existe en la magnitud que la discusión pública le atribuye, en un
aeropuerto, donde la vigilancia ya está normalizada, el efecto marginal
de añadir un sistema de conteo sería previsiblemente pequeño, aunque no
nulo, y la normalización misma formaría parte del fenómeno. Sostener el
punto con firmeza exigiría respaldo bibliográfico que este informe
todavía no tiene.

% ---- 4.4 ----
\subsection{Lectura ambiental del costo de cómputo}
\label{sec:impacto-ambiental}

El entrenamiento de modelos de visión es intensivo en cómputo y, por
tanto, en energía. El equipo no ha medido el consumo energético del
entrenamiento ni el de la inferencia, de modo que lo que sigue es
cualitativo y estructural; se indica en cada caso qué medición
permitiría convertirlo en cuantitativo.

Tres decisiones de diseño ya adoptadas o previstas tienen efecto directo
sobre la huella del entrenamiento. La primera es la reutilización de
modelos preentrenados: CSRNet parte de las primeras capas
convolucionales de VGG-16 ya entrenadas \cite{Li_2018_CVPR}, lo que
evita entrenar desde cero. La segunda es la eficiencia de etiquetado del
enfoque semi-supervisado de \cite{Lin_2025_CVPR}, sobre la que conviene
no sobrevender el argumento: reduce el esfuerzo humano de anotación, no
necesariamente el cómputo, puesto que el esquema maestro-estudiante
mantiene dos modelos en entrenamiento.

La tercera es la que admite una lectura ambiental más precisa, porque
los propios autores la cuantifican en prosa. P2R prescinde del algoritmo
húngaro, componente que los autores identifican expresamente como
costoso en tiempo dentro de la estrategia de emparejamiento punto a
punto. Sobre una imagen de 576 × 960 píxeles con 775 puntos anotados, el
cálculo de la pérdida punto a punto requiere 0,4307 segundos frente a
los 0,0064 segundos de P2R, cerca de sesenta y ocho veces más rápido
\cite{Lin_2025_CVPR}. La lectura ambiental de ese dato conviene hacerla
con cuidado, porque se presta a exageración: la cifra mide el costo de
una operación dentro de cada iteración de entrenamiento, no el consumo
total del ciclo, y el ahorro agregado depende del número de iteraciones
y de qué fracción del tiempo por iteración representaba esa operación.
Lo que sí puede afirmarse sin exceder el dato es que la elección de la
función de pérdida no es ambientalmente neutra, y que en un
entrenamiento con adaptación de dominio en sitio, que es el que la
brecha documentada obliga a contemplar, una diferencia de ese orden en
un componente ejecutado en cada iteración es un factor a considerar
junto con la exactitud.

A escala de la vida útil del sistema, el consumo de la inferencia
continua supera previsiblemente al del entrenamiento: este ocurre unas
pocas veces, aquella se ejecuta permanentemente sobre múltiples flujos
de video. Dos decisiones de diseño con motivación distinta convergen en
reducirlo. La frecuencia de muestreo es la primera: las aglomeraciones
se forman en escalas de minutos, de modo que procesar un fotograma cada
varios segundos en lugar de a velocidad completa reduce el cómputo en un
factor elevado sin pérdida operativa relevante, y reduce simultáneamente
el volumen de imágenes tratadas. La segunda es el procesamiento en el
borde, adoptado como medida M-01 por razones de protección de datos, que
evita transmitir video de forma continua hacia la nube con el ahorro de
transmisión asociado.

Este es el resultado analítico más útil de la subsección y conviene
presentarlo como tal ante el cliente: las decisiones que reducen la
exposición de datos personales (muestrear menos, procesar localmente, no
persistir imágenes, agregar la salida) son exactamente las mismas que
reducen el consumo de cómputo, de transmisión y de almacenamiento. Ambos
principios apuntan en la misma dirección, lo que refuerza el argumento
de privacidad con una justificación adicional que no depende del
cumplimiento normativo.

Quedan dos puntos abiertos. El primero es que convertir esta subsección
de cualitativa en cuantitativa depende de una sola medición: las horas
de GPU efectivamente empleadas durante el desarrollo. Es un registro
sencillo, pero no es reconstruible a posteriori: si no se lleva desde el
inicio del entrenamiento, la huella del desarrollo queda definitivamente
sin cuantificar. Por esa razón, este informe recomienda registrar las
horas de GPU desde el primer entrenamiento. El segundo es la huella de
hardware: si el proyecto requiere cámaras o servidores nuevos, su
fabricación y su disposición final integran el balance, y reutilizar la
infraestructura de CCTV existente es, además de la opción económicamente
decisiva, la ambientalmente preferible. Se recomienda, por último,
verificar si LAP cuenta con una política institucional de sostenibilidad
o con compromisos de reporte ambiental, porque, de existir, el análisis
debería alinearse con sus categorías de medición en lugar de proponer un
marco paralelo.

% ---- 4.5 ----
\subsection{Riesgos residuales y decisiones pendientes}
\label{sec:sintesis-impacto}
\label{sec:impacto-etico}

Integrando lo anterior, el perfil de riesgo del sistema admite una
síntesis en cuatro dimensiones. En la ética, el impacto positivo es
ofrecer una alternativa menos invasiva que la identificación biométrica
para la misma finalidad, y el riesgo principal es el deslizamiento hacia
identificación o seguimiento, mitigable mediante la declaración de
finalidades excluidas y las medidas M-02 y M-03. En la social, el
beneficio es la prevención de incidentes por aglomeración y la
asignación de recursos basada en evidencia, y los riesgos son la
distribución desigual de beneficio y carga, el uso laboral no declarado
y el exceso de confianza en el sistema. En la económica, el beneficio es
la eficiencia operativa y el valor de los incidentes evitados, y el
riesgo dominante es el costo de reconstruir el sistema si se detecta una
no conformidad después del despliegue, lo que hace de la privacidad
desde el diseño la opción económicamente racional además de la
obligatoria. En la ambiental, el beneficio es una huella menor que la de
alternativas de mayor granularidad, y el riesgo es el consumo acumulado
de la inferencia continua a lo largo de la vida útil.

Tres conclusiones transversales se desprenden del análisis. La primera
es que las decisiones de minimización operan simultáneamente en las
cuatro dimensiones: no persistir imágenes, agregar la salida y muestrear
a baja frecuencia reducen a la vez el riesgo ético, la exposición
social, el costo de almacenamiento y la huella ambiental. Son el núcleo
de la propuesta de este informe precisamente porque ningún objetivo
compite con otro.

La segunda es que el riesgo dominante es de gobernanza y no de
tecnología. El deslizamiento de finalidad y el uso no declarado no se
previenen con mejor ingeniería, sino con documentación, compromisos
contractuales y auditoría. Esto significa que la documentación de
cumplimiento no es accesoria al proyecto técnico: es la mitigación
principal del riesgo principal, y así debería exponerse ante el cliente.

La tercera es que varios impactos permanecen sin cuantificar por falta
de información del cliente y de mediciones propias, y se ha optado por
declararlo antes que por estimar sin base.

El riesgo de deslizamiento de finalidad merece un desarrollo propio,
porque es el principal y no proviene de un defecto técnico sino del
éxito del sistema. Una vez desplegada la infraestructura (cámaras
calibradas, canal de video, capacidad de inferencia, tablero de
visualización), el costo marginal de añadir capacidades es bajo:
incorporar reidentificación para medir tiempos de permanencia, o
reconocimiento facial para detectar personas de interés, es una decisión
de configuración y de presupuesto, no un rediseño. La barrera que impide
ese deslizamiento no es técnica sino documental y contractual. La
mitigación que este informe propone es documentar la finalidad de forma
restrictiva e incorporar al entregable una declaración de finalidades
excluidas: el sistema no realiza identificación de personas, no realiza
seguimiento individual de trayectorias, no realiza reconocimiento de
emociones y no alimenta decisiones automatizadas sobre pasajeros
individuales. Una lista de lo que el sistema \emph{no} hace es más
difícil de erosionar silenciosamente que una lista de lo que sí hace, y
su valor depende de que sea explícita: una finalidad excluida que solo
consta en el ánimo del equipo no restringe nada. A esa declaración
debería incorporarse expresamente la exclusión del uso laboral de los
datos, porque si los conteos por zona se cruzan con los turnos del
personal, el sistema se convierte de facto en una herramienta de
evaluación de desempeño, con implicancias laborales propias y con una
finalidad distinta de la declarada.

Los riesgos residuales que este informe no cierra son de tres clases, y
conviene no presentarlos como si fueran del mismo tipo. Los primeros son
ausencias de evidencia subsanables dentro del proyecto si se obtiene
acceso a material del terminal: la validación de exactitud en el dominio
de destino y la evaluación de sesgo demográfico. Los segundos son
decisiones del equipo que siguen abiertas y de las que depende buena
parte de este análisis: la arquitectura definitiva y la granularidad de
salida, la localización del procesamiento (\emph{on-premise} o nube
externa), el plazo concreto de retención y el registro de horas de GPU.
Los terceros son discrepancias entre documentos del propio proyecto, que
este informe deja registradas y expresamente no resuelve por su cuenta.
Son tres, y no todas tienen el mismo signo. El uso del término «personas
anonimizadas» para describir lo que la Ley 29733 califica como
disociación aparece tanto en la presentación preliminar del proyecto
como en la documentación del módulo de almacenamiento del prototipo, y
está tratado en la Sección 3.2. La medición de permanencia en zona
comercial como finalidad distinta de la estimación de aglomeraciones,
que exigiría base de legitimación propia, se trata en la Sección 3.2; el
esquema de datos del prototipo, que registra zona y marcas de primera y
última detección por persona, la hace técnicamente posible hoy. En
cambio, la preocupación sobre reidentificación mediante
\emph{embeddings} de apariencia \textbf{no se confirmó}: el prototipo
resuelve la fusión entre cámaras por asociación espacio-temporal y
reserva un descriptor de color no biométrico como mero desempate, de
modo que la decisión de diseño va en la dirección que exige el requisito
de no derivar de la imagen rasgos que permitan identificar a una
persona, ni siquiera de forma indirecta, según se detalla en la Sección
3.4.

A esas discrepancias se añade un riesgo residual de distinta clase,
porque no es una divergencia documental sino un incumplimiento
verificable: \textbf{la medida M-03 no se cumple en la implementación
actual}, que persiste coordenadas individuales enlazadas a un
identificador por persona en lugar de agregar por zona. Es el único
hallazgo de todo este análisis que el equipo puede cerrar por sí solo,
de inmediato y sin depender de LAP ni de acceso a datos del terminal, y
por eso debería encabezar la lista de acciones. Su cierre depende
únicamente del equipo, que puede agregar por zona antes de escribir o
fijar una retención muy corta sobre la tabla de posiciones; mientras no
lo haga, el cumplimiento de M-03 no puede afirmarse ante el cliente.

\clearpage
% =====================================================================
% METODOLOGIA -- SECCION A CARGO DE OTRO INTEGRANTE
% Responsable: no corresponde al frente etico-legal.
%
% NOTA PARA EL EQUIPO: las decisiones tecnicas que se documenten aqui son
% auditadas en la Seccion \ref{sec:auditoria} del marco etico-legal.
% Cualquier cambio de arquitectura, de dataset o de politica de retencion
% debe comunicarse al frente etico-legal, porque cambia la auditoria.
%
% Los marcadores \pendiente{PLACEHOLDER}{...} que este archivo tenia se
% retiraron del cuerpo. El encargo consta en PENDIENTES.md, seccion (b).
% =====================================================================
\section{Metodología}
\label{sec:metodologia}

\cajaaviso{Sección en elaboración}{El contenido de esta sección está a
cargo de otro integrante del equipo. El frente ético-legal la audita:
véase la Sección~\ref{sec:auditoria}. Dos decisiones que se documenten
aquí condicionan directamente esa auditoría (la arquitectura elegida y
el lugar donde se ejecuta la inferencia), por lo que conviene que se
fijen por escrito antes de cerrarla.}

\subsection{Arquitectura propuesta}
\label{sec:arquitectura}

Este apartado documentará la arquitectura del sistema de visión
computacional. La Sección~\ref{sec:sintesis-antecedentes} deja planteada
la elección entre los tres enfoques revisados y, sobre todo, la
granularidad de la salida, que es la decisión con mayor consecuencia para
el análisis de proporcionalidad.

\subsection{Datos y entrenamiento}
\label{sec:datos-entrenamiento}

Este apartado documentará los conjuntos de datos utilizados, la
estrategia de entrenamiento y la de adaptación de dominio. La
verificación de las licencias de los conjuntos públicos, que la
Sección~\ref{sec:licencias} deja pendiente, no se ha completado todavía
y depende de que el equipo documente aquí qué conjuntos emplea y bajo qué
licencia se distribuye cada uno.

\subsection{Métricas de evaluación}
\label{sec:metricas}

Este apartado justificará las métricas seleccionadas (MAE, MSE, nAP u
otras). La Sección~\ref{sec:ant-p2pnet} señala una restricción que
conviene tener presente al elegirlas: los compromisos de exactitud deben
formularse a nivel de zona agregada y no de posición individual.

\clearpage
% =====================================================================
% SECCION 9 -- RESULTADOS Y ANALISIS COMPARATIVO (esqueleto)
% Responsable: no corresponde al frente etico-legal.
%
% NOTA: la Seccion \ref{sec:ant-p2r} advierte que las cifras publicadas
% sobre datasets publicos no son extrapolables al terminal. Toda cifra
% que se reporte aqui debe indicar sobre que datos fue medida.
%
% Los marcadores \pendiente{PLACEHOLDER}{...} que este archivo tenia se
% retiraron del cuerpo. El encargo consta en PENDIENTES.md, seccion (b).
% =====================================================================
% 14/09/2026: el aviso "Seccion en elaboracion ... a cargo de otro
% integrante" se sustituyo por la frase neutra comun a todos los
% esqueletos. Reestructuracion al indice del modelo del curso.
\section{Resultados y Análisis Comparativo}
\label{sec:resultados}

Esta sección se completará en la versión final del informe.

\subsection{Resultados cuantitativos}
\label{sec:resultados-cuant}

Este apartado recogerá las tablas de error de conteo y de desempeño de
localización.

\subsection{Análisis cualitativo}
\label{sec:resultados-cual}

Este apartado recogerá casos de éxito y casos de fallo del modelo. Los
casos de fallo tienen aquí una función que excede lo técnico: la
Sección~\ref{sec:impacto-social} sostiene que los operadores deben
conocer los modos de fallo del sistema, y este apartado es la fuente
natural de esa información.

\clearpage
% =====================================================================
% SECCION 10 -- CONCLUSIONES Y RECOMENDACIONES (esqueleto)
% Responsable: no corresponde al frente etico-legal.
%
% NOTA: la sintesis de antecedentes (Seccion \ref{sec:sintesis-antecedentes})
% y la del analisis de impacto (Seccion \ref{sec:sintesis-impacto})
% alimentan esta seccion.
%
% Los marcadores \pendiente{PLACEHOLDER}{...} que este archivo tenia se
% retiraron del cuerpo. El encargo consta en PENDIENTES.md, seccion (b).
% =====================================================================
% 14/09/2026: el aviso "Seccion en elaboracion ... a cargo de otro
% integrante" se sustituyo por la frase neutra comun a todos los
% esqueletos. Reestructuracion al indice del modelo del curso.
\section{Conclusiones y Recomendaciones}
\label{sec:conclusiones}

Esta sección se completará en la versión final del informe.

Este apartado recogerá las conclusiones del proyecto y las líneas de
trabajo futuro. Dos resultados de la auditoría de la
Sección~\ref{sec:auditoria} deben reflejarse aquí sin atenuarlos: la
exactitud del sistema no está validada en el dominio de destino y el
desempeño desagregado por características demográficas no ha sido
evaluado. Ambos son subsanables dentro del proyecto, pero requieren
acceso a material del terminal, de modo que su cierre no depende
únicamente del equipo.

\clearpage

%------------------------ BIBLIOGRAFIA -------------------------------
\phantomsection
\addcontentsline{toc}{section}{Referencias}
\printbibliography[title={Referencias}]

\clearpage

%--------------------------- ANEXOS ----------------------------------
\appendix
% =====================================================================
% ANEXO A -- MODELO DE CARTEL INFORMATIVO
% Responsable: Fabian Moreno Ugarte
%
% El cartel se maqueta con \fbox + minipage + tabular, sin tcolorbox ni
% mdframed, para que compile en cualquier distribucion sin instalar
% paquetes adicionales. Desde 2026-09-14 va en negro, como todo el
% informe: el marco se conserva porque representa el borde fisico del
% cartel, no es un recuadro de aviso.
% =====================================================================
\section{Modelo de cartel informativo de videovigilancia}
\label{anx:cartel}

\subsection{Propósito y estatus de este modelo}

La Sección~\ref{sec:privacidad-diseno} estableció que el tratamiento de datos
personales mediante videovigilancia exige informar previamente a los
titulares mediante carteles visibles. Este anexo propone un modelo
concreto para el sistema del proyecto.

Se trata de una \textbf{propuesta de
trabajo}, no de un formato aprobado. Está construido sobre los elementos
informativos que la Directiva 01-2020-JUS/DGTAIPD exige, pero su
redacción final debe ser revisada por el área legal de LAP antes de su
colocación física.

Esa revisión tiene un contenido preciso: contrastar los campos del modelo
uno por uno contra el texto de la Directiva y contra el D.S.
016-2024-JUS, para confirmar que no falta ningún elemento informativo
exigido ni sobra alguno que la norma no requiera. Este informe propone el
conjunto de campos, no acredita que sea el exigible.

\subsection{Elementos que el modelo incorpora}

\begin{center}
\small
\begin{tabularx}{\linewidth}{@{}L{4.2cm}Y@{}}
\toprule
\textbf{Elemento} & \textbf{Razón de su inclusión} \\
\midrule
Identidad del responsable &
El titular debe saber ante quién ejercer sus derechos. \\
\addlinespace
Finalidad del tratamiento &
Núcleo del deber de información y límite frente al deslizamiento de
finalidad (Sec.~\ref{sec:impacto-etico}). \\
\addlinespace
Base de legitimación &
Permite al titular entender por qué no se le pidió consentimiento. \\
\addlinespace
Plazo de conservación &
Exigencia del principio de conservación limitada. \\
\addlinespace
Destinatarios &
Incluye la eventual entrega a autoridades conforme a las normas de
seguridad ciudadana (Sec.~\ref{sec:dl1218}). \\
\addlinespace
Canal para ejercer derechos &
Sin un canal indicado, el derecho es nominal. \\
\addlinespace
Declaración de lo que el sistema \emph{no} hace &
No es un elemento exigido por la norma. Se incorpora por decisión de
este frente: es la traducción visible del compromiso contra el
deslizamiento de finalidad y el elemento que distingue este cartel de
uno genérico de CCTV. \\
\bottomrule
\end{tabularx}
\end{center}

\subsection{Modelo propuesto}

En el modelo que sigue, los textos en \emph{cursiva} señalan los campos
que este informe no puede completar porque corresponden a datos que solo
LAP posee: su identificación registral, la base de legitimación que
decida invocar, sus canales de atención y el plazo de conservación que
fije. El resto del contenido es propuesta de este frente.

\vspace{0.5em}
\noindent
\fbox{%
\begin{minipage}{\dimexpr\linewidth-2\fboxsep-2\fboxrule\relax}
% Interlineado sencillo dentro del cartel: con el 1,5 del documento a
% 12 pt los filetes quedaban separados del texto por huecos muy grandes.
\setstretch{1}
\vspace{0.6em}
\begin{center}
{\Large\bfseries ZONA VIDEOVIGILADA}\\[0.25em]
{\large\bfseries SISTEMA DE ESTIMACIÓN DE AFORO}\\[0.2em]
{\small\itshape Este sistema NO identifica personas}
\end{center}
\noindent\rule{\linewidth}{0.8pt}
\vspace{0.2em}

\small
\setlength{\extrarowheight}{3pt}
\begin{tabularx}{\linewidth}{@{}L{3.1cm}Y@{}}
\textbf{Responsable} &
Lima Airport Partners S.R.L. (\emph{razón social registral completa,
RUC y domicilio fiscal, a completar por LAP}). \\

\textbf{Finalidad} &
Estimar el número de personas presentes por zona, con el fin de prevenir
aglomeraciones y gestionar el flujo de pasajeros en el terminal. \\

\textbf{Base legal} &
\emph{Base de legitimación que LAP invoque, con su fundamento normativo
exacto, a completar por su área legal.} \\

\textbf{Conservación} &
Las imágenes no se almacenan. El sistema conserva únicamente conteos
agregados por zona, durante el plazo fijado en su política de retención
(\emph{plazo a definir con LAP, medida M-05}). \\

\textbf{Destinatarios} &
Personal autorizado de LAP. Las imágenes del sistema de seguridad podrán
entregarse a las autoridades competentes cuando lo requieran conforme a
ley. \\

\textbf{Sus derechos} &
Puede ejercer sus derechos de acceso, rectificación, cancelación y
oposición escribiendo al canal de atención de LAP o acudiendo a su
oficina de atención al usuario (\emph{dirección de correo y ubicación
de la oficina, a completar por LAP}). \\

\textbf{Más información} &
\emph{Dirección web de la política de privacidad de LAP, a completar.} \\
\end{tabularx}

\vspace{0.4em}
\noindent\rule{\linewidth}{0.8pt}
\vspace{0.2em}

{\footnotesize\bfseries Este sistema no realiza reconocimiento facial, no
identifica a personas concretas, no realiza seguimiento individual de
trayectorias y no adopta decisiones automatizadas sobre los pasajeros.}

\vspace{0.6em}
\end{minipage}%
}
\vspace{0.8em}

\subsection{Criterios de colocación}

\begin{enumerate}[leftmargin=1.6em,itemsep=0.35em]
  \item \textbf{Anterioridad.} El cartel debe ser visible \emph{antes} de
        que la persona ingrese a la zona videovigilada, no dentro de
        ella. Informar a quien ya fue captado no cumple la finalidad de
        la norma.
  \item \textbf{Cobertura de todos los accesos.} Cada acceso a una zona
        cubierta requiere señalización. Un acceso sin cartel es un punto
        de incumplimiento con independencia de cuántos carteles haya en
        el resto del terminal.
  \item \textbf{Legibilidad.} Tamaño y contraste suficientes para ser
        leídos a la distancia de aproximación real, considerando que el
        pasajero circula con equipaje y no se detiene.
  \item \textbf{Idiomas.} Español e inglés como mínimo, dada la condición
        de terminal internacional. Si conviene añadir algún idioma más es
        algo que depende de la composición real del tráfico del
        aeropuerto, dato que debe pedirse a LAP y que este informe no
        supone.
  \item \textbf{Accesibilidad.} Altura y ubicación que permitan la
        lectura desde una silla de ruedas. El criterio se enuncia aquí en
        términos funcionales y no normativos: la normativa peruana de
        accesibilidad aplicable a señalización en establecimientos de uso
        público no se ha contrastado, de modo que las medidas concretas
        deben tomarse de ella y no de este anexo.
  \item \textbf{Registro de la colocación.} Documentar fotográficamente
        cada cartel instalado, con fecha y ubicación. Ante una
        fiscalización, el cumplimiento del deber de información debe
        poder acreditarse, y la fotografía fechada es la evidencia más
        simple de producir.
\end{enumerate}

\subsection{Observación final sobre este anexo}

Este cartel es el punto donde el análisis normativo de este informe se
vuelve visible para el pasajero. Todo lo demás (la arquitectura, las
medidas de privacidad desde el diseño, la auditoría) ocurre en
documentos que el titular de los datos nunca verá. La frase
\enquote{este sistema no identifica personas} solo es defendible si las
medidas M-02, M-03 y M-04 de la Sección~\ref{sec:privacidad-diseno}
están efectivamente implementadas.

Colocar un cartel que declare más de lo que el
sistema cumple es peor que no colocarlo: convierte una deficiencia
técnica en una declaración inexacta hacia el titular. La redacción final
del cartel debe fijarse \emph{después} de que la arquitectura esté
congelada, y debe revisarse si la arquitectura cambia.


\end{document}
